\section{UIOrthoLoRA - Unit Invariant Flexibility for Structural Adaptation}
\label{sec:uiortholora}

While OrthoLoRA addresses the fixed-basis bottleneck by allowing rotations within
the pretrained low-spectrum subspace, it still relies on a strong form of the
pretraining-finetuning paradigm: the assumption that task-relevant adaptation
can be expressed entirely within the singular geometry inherited from
pretraining.

To relax this assumption while retaining structural preservation, we introduce
trainable diagonal scalers $E$ and $D$, inspired by the unit-invariant SVD
perspective of \citet{uhlmann2018generalized}. These scalers apply lightweight
diagonal transformations to the output and input coordinates, respectively:
\[
\Delta W^{ui}
=
\underline E U_r \underline R_U \underline \Sigma'_r \underline R_V^\top V_r^\top \underline D .
\]
Thus, $R_U$ and $R_V$ provide directional flexibility within the pretrained
low-spectrum subspace, while $E$ and $D$ provide unit-level coordinate
adaptation around this subspace.

This diagonal scaling eases the strict reliance on the pretrained singular basis. Although the update remains tied to the UI-SVD parameterization, $E$ and $D$ can induce controlled mixing between pretrained singular coordinates. This gives the adapter additional flexibility to express task-specific perturbations that are not perfectly aligned with the original tail singular directions. The linear version, UILinLoRA, disables the orthogonal rotations ($R_U = R_V = I$), and keeps diagonal scalers in place.

\subsection{Scaler-Induced Mixing in Pretrained Singular Coordinates}
Throughout this section, \(U\), \(V\), and their partitions \(U_R, U_r, V_R, V_r\) refer to the pretrained singular bases. The diagonal scalers $E$ and $D$ relax the strict confinement assumption of
OrthoLoRA. To analyze their effect, define the rotated tail bases $\bar{U}_r := U_r R_U, \bar{V}_r := V_r R_V .$
Since $R_U$ and $R_V$ are orthogonal, $\bar{U}_r$ and $\bar{V}_r$ span the same
tail subspaces as $U_r$ and $V_r$, respectively. Therefore, without the diagonal
scalers, OrthoLoRA remains exactly confined to the pretrained tail singular
block.

With unit-invariant scaling, the update becomes
\[
\Delta W^{ui}
=
\underline E \bar{U}_r \underline \Sigma'_r \bar{V}_r^\top  \underline D .
\]
The scalers $E$ and $D$ can change how the rotated tail directions are expressed
in the original feature coordinates. This gives the adapter additional
flexibility, but it may also introduce coupling with the preserved leading
singular structure.

To quantify this effect, we write the perturbation in the pretrained singular coordinates: \\ $C := U^\top \Delta W^{ui} V$.
Partitioning $C$ according to the leading-tail decomposition gives
\[
C
=
\begin{bmatrix}
C_{11} & C_{12} \\
C_{21} & C_{22}
\end{bmatrix}
=
\begin{bmatrix}
U_R^\top E \bar{U}_r \Sigma'_r \bar{V}_r^\top D V_R
&
U_R^\top E \bar{U}_r \Sigma'_r \bar{V}_r^\top D V_r
\\[0.5em]
U_r^\top E \bar{U}_r \Sigma'_r \bar{V}_r^\top D V_R
&
U_r^\top E \bar{U}_r \Sigma'_r \bar{V}_r^\top D V_r
\end{bmatrix}.
\]

The blocks have the following interpretation:
\[
C_{11}: \text{direct leakage into the preserved leading block},
\]
\[
C_{12}, C_{21}: \text{cross-coupling between leading and tail coordinates},
\]
\[
C_{22}: \text{adaptation internal to the tail block}.
\]

We define the left and right scaler-induced mixing terms as
\[
M_E := U_R^\top E \bar{U}_r,
\quad
M_D := V_R^\top D \bar{V}_r
\text{\quad with magnitudes \quad}
\mu_E := \|M_E\|_2,
\quad
\nu_D := \|M_D\|_2 .
\]

Here, $\mu_E$ measures how much the scaled left tail basis $E\bar{U}_r$ enters
the preserved left subspace, and $\nu_D$ measures how much the scaled right tail
basis $D\bar{V}_r$ enters the preserved right subspace.

\noindent \textbf{Proposition 2} (Scaler-Induced Mixing Bounds).
\textit{
Let
\[
\Delta W^{ui}
=
E \bar{U}_r \Sigma'_r \bar{V}_r^\top D, \quad
C := U^\top \Delta W^{ui} V
=
\begin{bmatrix}
C_{11} & C_{12} \\
C_{21} & C_{22}
\end{bmatrix}.
\]
Then the perturbation blocks satisfy
}
\[
\|C_{11}\|_2
\leq
\mu_E \|\Sigma'_r\|_2 \nu_D,
\]
\[
\|C_{12}\|_2
\leq
\mu_E \|\Sigma'_r\|_2 \|D\|_2,
\]
\[
\|C_{21}\|_2
\leq
\|E\|_2 \|\Sigma'_r\|_2 \nu_D,
\]
\[
\|C_{22}\|_2
\leq
\|E\|_2 \|\Sigma'_r\|_2 \|D\|_2.
\]


\textit{Proof is provided in Appendix~\ref{app:proof-proposition-2-scaler-induces-mixing-bounds}.}

\noindent \textbf{Corollary 1} (Exact Tail Confinement as a Special Case).
\textit{
If $\mu_E = 0$ and $\nu_D = 0$, then:
\[
C_{11}=0,\qquad C_{12}=0,\qquad C_{21}=0 
\implies
C
=
\begin{bmatrix}
0 & 0 \\
0 & C_{22}
\end{bmatrix}.
\]
}

This means that the update is confined entirely to the tail block in pretrained
singular coordinates. In particular, pure OrthoLoRA corresponds to the special
case $E=I$ and $D=I$, for which
\[
U_R^\top \bar{U}_r = 0,
\qquad
V_R^\top \bar{V}_r = 0,
\]
so the update remains strictly orthogonal to the preserved leading singular
structure.

\noindent \textbf{Remark} (Interpretation).
The diagonal scalers $E$ and $D$ should be viewed as a controlled relaxation of
strict tail confinement. When $E=I$ and $D=I$, UIOrthoLoRA adapts only within the pretrained tail subspace. When $E$ and $D$ are learned, the adapter can reweight
the input and output coordinates, allowing the effective adapted directions
$E\bar{U}_r$ and $D\bar{V}_r$ to deviate from the original tail coordinates.

This deviation is not unstructured: it is exactly measured by
\[
\mu_E = \|U_R^\top E\bar{U}_r\|_2,
\qquad
\nu_D = \|V_R^\top D\bar{V}_r\|_2.
\]
Thus, $E$ and $D$ provide additional expressive capacity, while the quantities
$\mu_E$ and $\nu_D$ quantify the extent to which this capacity interacts with
the preserved dominant singular structure.


\noindent \textbf{Theorem 2} (Approximate Tail Preservation under Small Mixing).
Let
\[
C := U^\top \Delta W^{ui} V
=
\begin{bmatrix}
C_{11} & C_{12} \\
C_{21} & C_{22}
\end{bmatrix},
\quad
C_{\mathrm{tail}}
=
\begin{bmatrix}
0 & 0 \\
0 & C_{22}
\end{bmatrix}
\]
be the UIOrthoLoRA update written in the pretrained singular coordinates, and its pure tail-block component. Then
\[
\|C-C_{\mathrm{tail}}\|_2
\leq
\|\Sigma'_r\|_2
\sqrt{
(\mu_E\nu_D)^2
+
(\mu_E\|D\|_2)^2
+
(\|E\|_2\nu_D)^2
}.
\]
Consequently, if $\|E\|_2$, $\|D\|_2$, and $\|\Sigma'_r\|_2$ remain bounded, and $\mu_E \to 0,
\nu_D \to 0 \\
\implies
\|C-C_{\mathrm{tail}}\|_2 \to 0.
$

Thus, small scaler-induced mixing implies that the UIOrthoLoRA update remains
approximately confined to the pretrained tail singular block, even though the
diagonal scalers $E$ and $D$ relax strict tail confinement.

\textit{Proof is provided in Appendix~\ref{app:proof-thm-2-approximate-tail-preservation}.}

Appendix~\ref{app:soft_mixing_ablation} provides a diagnostic ablation showing
that these mixing quantities can also be used as soft regularizers: penalizing
both left- and right-side mixing reduces measured off-tail leakage and
leading-subspace drift on 5 GLUE tasks while maintaining task performance.


\noindent \textbf{Remark} (Gap-Dependent Dominant Subspace Stability).
The mixing bounds above control how much the UIOrthoLoRA update enters the
preserved leading block in pretrained singular coordinates. A complementary
view is obtained from standard singular subspace perturbation theory. Let
\[
\gamma := \sigma_R(W^{pre}) - \sigma_{R+1}(W^{pre}) > 0
\]
denote the spectral gap between the preserved leading singular components and
the remaining tail components. If $\|\Delta W^{ui}\|_2 < \gamma$, then the leading left and right singular subspaces of $W^{pre}+\Delta W^{ui}$
remain close to those of \(W^{pre}\), with subspace movement scaling on the
order of
\[
\mathcal{O}\left(\frac{\|\Delta W^{ui}\|_2}{\gamma}\right).
\]
Since $\|\Delta W^{ui}\|_2
\leq
\|E\|_2\|\Sigma'_r\|_2\|D\|_2,
$
the overall scaled update magnitude controls global dominant-subspace drift,
while the sharper block-level quantity $\|C_{11}\|_2
\leq
\mu_E\|\Sigma'_r\|_2\nu_D
$
controls direct leakage into the preserved leading singular coordinates.

